\chapter{Ising model with the Metropolis algorithm}

\section{Simulations at different temperatures}

\begin{sourcefiles}
  \item[Source code] src/051\_ising\_metropolis.c
  \item[Analysis script] src/051\_ising\_metropolis.R
  \item[Executable] exe/051\_ising\_metropolis
\end{sourcefiles}

For this and the following exercise, I implemented the Metropolis algorithm with
Glauber dynamics and sequential updating to simulate the two-dimensional Ising
model. The program outputs the total energy and magnetization after each sweep
of $N = L^{2}$ local spin flips.

In order to correctly estimate the observables of interest, after discarding an
appropriate number of equilibration time steps, I calculated the autocorrelation
time of each time series. Given an observable $\mathcal{O}$, we estimate its
expectation value $\expval{\mathcal{O}}$ using the Monte Carlo average
\begin{equation}
  \bar{\mathcal{O}} = \frac{1}{T} \sum_{t=1}^{T} \mathcal{O}_{t},
\end{equation}
where $T$ is the length of the time series after equilibration. The sample
variance of the estimator can be derived from the expression
\begin{align}
  \sigma_{\bar{\mathcal{O}}}^{2} &= \expval{\bar{\mathcal{O}}^{2}} -
  \expval{\bar{\mathcal{O}}}^{2} = \frac{1}{T(T - 1)} \sum_{s,t = 1}^{T}
  \expval{\mathcal{O}_{s} \mathcal{O}_{t}} - \frac{1}{T(T - 1)} \sum_{s,t =
  1}^{T} \expval{\mathcal{O}_{s}} \expval{\mathcal{O}_{t}}. \label{eq:var-obar}
  \\
\intertext{If we collect the $s = t$ and $s \neq t$ terms, we get}
  \sigma_{\bar{\mathcal{O}}}^{2} &= \frac{1}{T(T - 1)} \sum_{t=1}^{T}
  \left(\expval{\mathcal{O}_{t}^{2}} - \expval{\mathcal{O}_{t}}^{2}\right)
  + \frac{1}{T(T - 1)} \sum_{s \neq t}^{T} \left(\expval{\mathcal{O}_{s}
  \mathcal{O}_{t}} - \expval{\mathcal{O}_{s}} \expval{\mathcal{O}_{t}}\right).
\end{align}
The first term is simply the variance of the time series,
$\sigma_{\mathcal{O}}^{2}$, divided by $T - 1$. For the second term, we can
rewrite it as
\begin{align}
  \sum_{s \neq t}^{T} \left(\expval{\mathcal{O}_{s} \mathcal{O}_{t}} -
  \expval{\mathcal{O}_{s}} \expval{\mathcal{O}_{t}}\right) &= 2 \sum_{s = 1}^{T
  - 1} \sum_{t = 1}^{T - s} \left(\expval{\mathcal{O}_{s} \mathcal{O}_{t}} -
  \expval{\mathcal{O}_{s}} \expval{\mathcal{O}_{t}}\right), \\
\intertext{and then, since we are at equilibrium, we assume translational
invariance to group together all the pairs at distance $t$, of which there are
$T - t$:}
  \sum_{s \neq t}^{T} \left(\expval{\mathcal{O}_{s} \mathcal{O}_{t}} -
  \expval{\mathcal{O}_{s}} \expval{\mathcal{O}_{t}}\right) &=
  2 \sum_{t=1}^{T} \left(\expval{\mathcal{O}_{s} \mathcal{O}_{s+t}} -
  \expval{\mathcal{O}_{s}} \expval{\mathcal{O}_{s+t}}\right) (T - t),
\end{align}
where the $t = T$ cancels out. Inserting this result back into
\cref{eq:var-obar}, we can express:
\begin{equation}
  \sigma_{\bar{\mathcal{O}}}^{2} = \frac{\sigma_{\mathcal{O}}^{2}}{T} +
  \frac{2}{T} \sum_{t=1}^{T} \left(\expval{\mathcal{O}_{s}
  \mathcal{O}_{s+t}} - \expval{\mathcal{O}_{s}}
  \expval{\mathcal{O}_{s+t}}\right) \left(\frac{T - t}{T - 1}\right) 
  = \frac{\sigma_{\mathcal{O}}^{2}}{T} (1 + 2 \tau_{\mathcal{O}}),
\end{equation}
where we define the \emph{integrated autocorrelation time} as
\begin{equation}
  \tau_{\mathcal{O}} \coloneq \sum_{t=1}^{T} \left(\frac{T-t}{T-1}\right)
  C_{\mathcal{O}}(t),
  \label{eq:tau-def}
\end{equation}
with $C_{\mathcal{O}}$ being the \emph{autocorrelation function} defined as
\begin{equation}
  C_{\mathcal{O}}(t) \coloneq
  \frac{
    \expval{\mathcal{O}_{s} \mathcal{O}_{s + t}}
    - \expval{\mathcal{O}_{s}} \expval{\mathcal{O}_{s+t}}
  }{
    \expval{\mathcal{O}_{s}^{2}} - \expval{\mathcal{O}_{s}}^{2}
  }
\end{equation}

Although I used \cref{eq:tau-def} in my codes, one could simplify it knowing
that the autocorrelation function decays exponentially for $t \to \infty$:
\begin{equation}
  C_{\mathcal{O}}(t) \sim e^{-t / \tau_{\mathcal{O}}'},
\end{equation}
so that we can approximate the integrated correlation time with
\begin{equation}
  \tau_{\mathcal{O}}' \coloneq \sum_{t=1}^{T} C_{\mathcal{O}}(t) \approx
  \tau_{\mathcal{O}}.
\end{equation}
